\section{QARMAs}

\begin{enumerate}
  \item What rules of engagement should apply when suspected vessels refuse to comply with boarding requests in international waters, given that vessels enjoy freedom of navigation under UNCLOS?
  \item How should NATO deal with the evidentiary challenges of proving intent, given that anchor-dragging incidents are designed to appear as plausible accidents rather than sabotage, as demonstrated by the dismissal of the Eagle S crew case in Finland?
  \item How should NATO enhance its military surveillance and patrol presence in the Baltic Sea beyond the current Baltic Sentry operation, and what additional assets (frigates, drones, maritime patrol aircraft) should be deployed?
  \item How can the Alliance address the capability gap of frontline Baltic states (Estonia, Latvia, Lithuania), which are the most exposed but have historically prioritized land forces over naval capacity?
  \item What role should AI-driven surveillance systems, such as the UK's Nordic Warden system, play in tracking and flagging shadow fleet vessels across the region?
  \item What role should public-private cooperation play, following models like Lithuania's agreement between its armed forces and grid operator Litgrid?
\end{enumerate}
